\documentclass[9pt,a4paper]{extarticle}

% ======================
% Codificação, idioma e layout
% ======================
\usepackage[T1]{fontenc}
\usepackage[utf8]{inputenc}
\usepackage[brazil]{babel}
\usepackage{lmodern}
\usepackage{geometry}
\geometry{margin=2cm}
\usepackage{parskip}
\usepackage{microtype}

% ======================
% Matemática e teoremas
% ======================
\usepackage{amsmath, amssymb, amsfonts, bm}
\newtheorem{definition}{Definição}

% ======================
% Listas
% ======================
\usepackage{enumitem}
\setlist{nosep}

% ======================
% Macros
% ======================
\newcommand{\sigmoid}{\sigma}
\newcommand{\bw}{\bm w}
\newcommand{\bx}{\bm x}
\newcommand{\by}{\bm y}


% ======================
% Documento
% ======================
\begin{document}

\title{\vspace{-1em}Prova 4\\
Aprendizado de Máquina}
\date{}
\maketitle

\vspace{-2cm}

\section*{Instruções gerais}
\begin{enumerate}
\item Todo o material de consulta deverá ser entregue juntamente com a prova, ao final da avaliação.
\item O único material impresso permitido são as listas de exercícios disponibilizadas no repositório oficial da disciplina.
\item Com exceção do material mencionado no item 2, todo o restante do material de consulta deverá ser escrito de próprio punho.
\item Não será permitido o uso de dispositivos eletrônicos, tais como celulares, tablets, notebooks, smartwatches ou similares, durante a prova.
\item Não é permitido o compartilhamento de material de consulta entre estudantes durante a realização da prova.
\end{enumerate}

% ==================================================
\section*{Questões}
% ==================================================

\begin{enumerate}[label=\textbf{\arabic*)}, itemsep=1.3em]
\item Para cada aplicação abaixo indique o tipo de aprendizado de máquina mais adequado (\emph{supervisionado}, \emph{não supervisionado} ou \emph{por reforço}) e justifique a sua indicação:

\begin{enumerate}[label=(\alph*)]
    \item Uma operadora de telecomunicações deseja identificar automaticamente grupos de clientes com perfis de consumo semelhantes, sem possuir rótulos ou categorias previamente definidos.

    \item Uma empresa de streaming deseja prever a nota que um usuário atribuiria a um filme, com base no histórico de avaliações de usuários para diversos filmes.
    
    \item Uma empresa de logística deseja controlar automaticamente a rota de um veículo autônomo em um armazém, de forma a minimizar o tempo de entrega e evitar colisões, aprendendo a partir da interação contínua com o ambiente.
\end{enumerate}

\item Para cada aplicação abaixo, indique o algoritmo de aprendizado supervisionado mais adequado 
(\emph{regressão linear}, \emph{regressão logística} ou \emph{redes neurais artificiais}) 
e justifique a sua escolha. (\textbf{Obs.:} Mais de uma alternativa pode ser considerada adequada em cada aplicação. A avaliação levará em conta a capacidade de justificar tecnicamente a escolha dos algoritmos.):

\begin{enumerate}[label=(\alph*)]
    \item Uma empresa deseja prever o tempo (em minutos) que um usuário permanecerá em uma plataforma digital, a partir de variáveis como horário de acesso, tipo de dispositivo, histórico recente de uso e perfil do usuário.
    
    \item Uma instituição financeira deseja estimar a probabilidade de inadimplência de um cliente, com base em renda, histórico de crédito, número de atrasos anteriores e valor do empréstimo solicitado.
    
    \item Um sistema de visão computacional deve identificar se uma imagem contém ou não um determinado objeto, a partir de descritores extraídos automaticamente das imagens.
\end{enumerate}


\item O que é e para o que serve uma função de $Loss$?

\item Qual a relação entre aprendizado de máquina e otimização?

\item Por que estudamos como obter o gradiente das funções de loss? O gradientes são calculados em função de quais variáveis num problema de aprendizagem de máquina?

\item Para cada uma das expressões abaixo, indique:  i) o que elas representam, ii) o que cada variável significa, e iii) como são utilizadas no aprendizado de máquina. 
\begin{enumerate}[label=(\alph*)]
    \item \[ \min_{\bm w}\; \lVert X\bm w - \bm y \rVert^2 \]
    \item \[ J(\bw)= -\sum_{i=1}^n \Big(y_i\log p_i + (1-y_i)\log(1-p_i)\Big) \]
    \item \[ f(x) = \max(0, x) \]

\end{enumerate}


\item Considere os dois cenários a seguir, relacionados ao desenvolvimento e à avaliação de modelos de aprendizado de máquina. Para cada cenário (i) identifique se ocorre o vazamento de dados; (ii) em caso positivo, explique por que esse procedimento compromete a estimativa de desempenho do modelo; e (iii) descreva como o processo deveria ser modificado para evitar o vazamento.

\begin{enumerate}[label=(\alph*)]
    \item Em um problema de classificação, um pesquisador realiza a seleção de atributos utilizando todo o conjunto de dados disponível (por exemplo, escolhendo as variáveis mais correlacionadas com o rótulo). Em seguida, o conjunto de dados reduzido é dividido em treino e teste, e um classificador é treinado e avaliado.

    \item Em um problema de regressão, os dados são inicialmente normalizados por meio do cálculo da média e do desvio padrão usando todas as amostras. Após esse pré-processamento, o conjunto é dividido em treino e teste, e o modelo é ajustado e avaliado.
\end{enumerate}


\item Compare os papéis dos conjuntos de treinamento, validação e teste em um fluxo adequado de desenvolvimento de modelos preditivos. Quais são as consequências metodológicas de utilizar o conjunto de teste para tomada de decisões de modelagem?

\item Considere um modelo cujo ajuste de hiperparâmetros foi realizado por validação cruzada. Descreva o procedimento correto para obter uma estimativa confiável do desempenho do modelo em dados não observados.

\item Explique como a análise do desempenho de um modelo em função do número de iterações de treinamento pode ser usada para diagnosticar problemas de subajuste (uderfitting) e sobreajuste (overfitting). Descreva os comportamentos típicos esperados nessas situações.

 


\end{enumerate}

\end{document}
